\chapter{Swollen Lymph Nodes}

\section*{Definition}
Lymphadenopathy refers to enlargement of lymph nodes, typically exceeding 1.0~cm in diameter for cervical, axillary, or inguinal nodes, or 0.5~cm for epitrochlear nodes. Localized lymphadenopathy involves a single region while generalized lymphadenopathy involves two or more non-contiguous regions and raises suspicion for systemic disease.

\section*{What Happens in Your Body}
Lymph nodes enlarge through three principal mechanisms: reactive hyperplasia from antigenic stimulation (increased lymphocytes and macrophages proliferating within germinal centers), infiltration by malignant cells (lymphoma, metastatic carcinoma, leukemia), or accumulation of inflammatory cells and organisms in infectious lymphadenitis. Antigen-presenting cells in the node activate T- and B-lymphocyte proliferation, producing the palpable enlargement that resolves as the immune response subsides.

\section*{Causes}
\begin{itemize}
  \item Infectious causes: Viral URI (EBV, CMV, HIV, influenza, adenovirus), Bacterial pharyngitis (Group A Strep), Tuberculosis, Atypical mycobacteria, Cat-scratch disease (Bartonella henselae), Toxoplasmosis, Syphilis, Lymphogranuloma venereum, Fungal infections (histoplasmosis, coccidioidomycosis)
  \item Malignant causes: Hodgkin lymphoma, Non-Hodgkin lymphoma, Chronic lymphocytic leukemia, Acute leukemias, Metastatic carcinoma (head and neck, breast, lung, melanoma)
  \item Autoimmune causes: Systemic lupus erythematosus, Rheumatoid arthritis, Sarcoidosis, Still disease, IgG4-related disease
  \item Other causes: Drug hypersensitivity (phenytoin, allopurinol), Storage diseases (Gaucher, Niemann-Pick), Post-vaccination, Kawasaki disease, Kikuchi disease
\end{itemize}

\section*{When to See a Doctor}
See your doctor if:
\begin{itemize}
  \item Node $>$ 2.5~cm in diameter, hard, fixed, or matted (suspicious for malignancy)
  \item Supraclavicular or scalene node involvement (Virchow node suggests GI malignancy)
  \item Generalized lymphadenopathy without obvious infectious cause
  \item Persistent enlargement beyond 4 weeks despite observation
  \item Associated with unexplained fever, night sweats, weight loss (B symptoms)
  \item Overlying redness, warmth, fluctuance suggesting suppurative lymphadenitis
\end{itemize}

\section*{Related Symptoms}
\begin{itemize}
  \item Fever
  \item Night sweats
  \item Weight loss
  \item Fatigue
  \item Splenomegaly
  \item Hepatomegaly
  \item Pharyngitis
  \item joint pain
  \item Rash
  \item itching
\end{itemize}
\textbf{Important:} Imaging (ultrasound, CT) and biopsy (excisional preferred over FNA for lymphoma diagnosis) guided by node characteristics, location, and clinical context.

\section*{Self-Care / Home Management}
\begin{itemize}
  \item Observe small ($<$ 1~cm), soft, mobile nodes for 2--4 weeks, especially with concurrent URI symptoms.
  \item Warm compresses may provide comfort for tender nodes.
  \item Acetaminophen or NSAIDs for associated pain or fever.
  \item Avoid repeatedly palpating nodes as this may increase inflammation.
\end{itemize}

\section*{How Your Doctor Will Evaluate You}
\begin{itemize}
  \item History focuses on symptom duration, growth pattern, B symptoms (fever, night sweats, weight loss), exposures (pets, travel, sick contacts), and medication review.
  \item Physical examination documents node size, consistency, mobility, tenderness, and location; assess for splenomegaly and hepatomegaly.
  \item Initial labs include CBC with differential, peripheral smear, ESR, CRP, LDH, uric acid; targeted serology for EBV, CMV, HIV, VDRL, or PPD/IGRA.
  \item Imaging includes ultrasound for superficial nodes or CT chest/abdomen/pelvis for deep nodes.
  \item Excisional biopsy of the largest, most suspicious node for histology, flow cytometry, and cytogenetics when malignancy suspected.
\end{itemize}

\section*{Treatment Options}
\begin{itemize}
  \item Reactive lymphadenopathy from viral illness is managed with supportive care and resolves spontaneously.
  \item Bacterial lymphadenitis requires antibiotics directed at the suspected organism (amoxicillin-clavulanate, clindamycin, or TMP-SMX).
  \item Suppurative nodes may require incision and drainage with culture.
  \item Tuberculous lymphadenitis treated with standard four-drug anti-TB regimen.
  \item Malignant lymphadenopathy requires oncology referral for chemotherapy, radiation, or targeted therapy per histologic type.
\end{itemize}

\section*{References}
\begin{thebibliography}{99}
\bibitem{swollen_lymph_nodes:ref1} Ferrer RL. "Lymphadenopathy: Differential diagnosis and evaluation." \textit{American Family Physician}, 1998;58(6):1313--1320.
\bibitem{swollen_lymph_nodes:ref2} Bazemore AW, Smucker DR. "Lymphadenopathy and malignancy." \textit{American Family Physician}, 2002;66(11):2103--2110.
\bibitem{swollen_lymph_nodes:ref3} Heilen M, Schmid M, Grasshoff C, et al. "Cervical lymphadenopathy in children: A diagnostic algorithm." \textit{Pediatric Reports}, 2021;13(2):280--289.
\bibitem{swollen_lymph_nodes:ref4} Pangalis GA, Vassilakopoulos TP, Boussiotis VA, et al. "Clinical approach to lymphadenopathy." \textit{Seminars in Oncology}, 1993;20(6):570--582.
\bibitem{swollen_lymph_nodes:ref5} Habermann TM, Steensma DP. "Lymphadenopathy." \textit{Mayo Clinic Proceedings}, 2000;75(7):723--732.
\bibitem{swollen_lymph_nodes:ref6} Schrager L, Friedland GH. "Lymphadenopathy." In: \textit{Harrison's Principles of Internal Medicine}, 21st ed. McGraw-Hill, 2022.
\end{thebibliography}
